% ---------------------------------------------------------------- title page
\begin{titlepage}
\begin{center}
\vspace*{0.2cm}
{\LARGE\bfseries \projecttitle\par}
\vspace{1.2cm}
{\large A Project Report\\[0.2em] submitted in partial fulfilment of the requirements\\ for the award of the degree of\par}
\vspace{0.6cm}
{\Large\bfseries \degree\\[0.2em] in\\[0.2em] \branch\par}
\vspace{1.0cm}
{\large Submitted by\par}
\vspace{0.2cm}
{\large\bfseries \studentname\par}
{\large \rollno\par}
\vspace{0.9cm}
{\large Under the supervision of\par}
\vspace{0.2cm}
{\large\bfseries \supervisor\par}
{\large \supervisordesig\par}
\vfill
% Add the DTU logo here if your department requires it:
% \includegraphics[width=3cm]{dtu_logo.png}\\[0.6cm]
{\large\bfseries \department\par}
{\large\bfseries \university\par}
{\large \universityaddr\par}
\vspace{0.4cm}
{\large \submissionmonth\par}
\end{center}
\end{titlepage}

% --------------------------------------------------------- declaration
\chapter*{Candidate's Declaration}
\addcontentsline{toc}{chapter}{Candidate's Declaration}
I, \studentname{} (\rollno), a student of \degree{} in \branch{}, hereby declare that the
project report titled \emph{``\projecttitleplain''}, submitted to the \department{},
\university{}, in partial fulfilment of the requirements for the award of the degree of
\degree{}, is an original record of my work carried out under the supervision of \supervisor{}.
This work has not been submitted, in part or in full, for the award of any other degree or
diploma at this or any other institution.

\vspace{2.5cm}
\noindent Place: Delhi \hfill \studentname\\
Date: \fillme{Date} \hfill \rollno

% --------------------------------------------------------- certificate
\chapter*{Certificate}
\addcontentsline{toc}{chapter}{Certificate}
I hereby certify that the project report titled \emph{``\projecttitleplain''}, submitted by
\studentname{} (\rollno), \department{}, \university{}, in partial fulfilment of the
requirements for the award of the degree of \degree{} in \branch{}, is a record of the
project work carried out by the student under my supervision. To the best of my knowledge
this work has not been submitted, in part or in full, for any degree or diploma at this
university or elsewhere.

\vspace{2.5cm}
\noindent Place: Delhi \hfill \supervisor\\
Date: \fillme{Date} \hfill \supervisordesig\\
\mbox{} \hfill \department\\
\mbox{} \hfill \university

% --------------------------------------------------------- acknowledgement
\chapter*{Acknowledgement}
\addcontentsline{toc}{chapter}{Acknowledgement}
I am grateful to my supervisor, \supervisor{}, for guidance and feedback throughout this
project, and to the faculty of the \department{}, \university{}, for the coursework in
probability, stochastic processes, optimisation and computing on which this work rests.
\fillme{Add or edit personal acknowledgements: family, friends, lab staff, etc.}

\vspace{2cm}
\noindent\hfill \studentname

% --------------------------------------------------------- abstract
\chapter*{Abstract}
\addcontentsline{toc}{chapter}{Abstract}
Internet-facing SSH servers are attacked continuously, yet conventional defences only block
attackers and learn nothing about what they would have done after logging in. A honeypot
inverts this: it admits the attacker into a decoy system and records their behaviour.
Existing SSH honeypots respond to every command with the same fixed emulation. This project
asks whether a honeypot that \emph{chooses} how to respond to each command can capture more
attacker behaviour than a static one.

We design and build \textbf{CruxTrap}, a medium-interaction SSH honeypot that emulates an
Ubuntu server inside a per-session, in-memory filesystem and never executes attacker input.
Each command is classified into one of ten intent categories. A tabular reinforcement-learning
agent observes the pair (category, session phase) and selects one of five responses:
\act{allow}, \act{delay}, \act{fake} (serve decoy credentials), \act{tarpit} or \act{block}.
Events flow through Redis Streams to a pipeline that enriches them with GeoIP data and stores
them in PostgreSQL, and a SOC-style dashboard supports investigation and session replay.
The five services are containerised with defence-in-depth hardening, and the system's
behaviour under component failure is measured.

The interaction is formalised as a Markov decision process whose reward measures captured
intelligence. Because the attacker's type is hidden, the honeypot faces a partially
observable problem. To evaluate policies rigorously we build a generative model of three
attacker types and compute each policy's expected performance \emph{exactly}, by solving
the Bellman equations on a 330-state augmented chain. We also compute an oracle upper bound
by value iteration and the best policy available to the honeypot's own state representation.
Over 20 independent runs, Q-learning outperforms a static honeypot (expected discounted return
$5.686 \pm 0.051$ versus $5.542$, $p = 5.8 \times 10^{-6}$) and recovers 60\% of the achievable
improvement. It does so by learning, without supervision, to serve decoy credentials to attackers
who are searching for credentials. With the conventional constant learning rate, however, the
agent performed \emph{worse} than the static baseline. Stochastic-approximation theory explains
why, and an ablation confirms that a Robbins--Monro step-size schedule fixes it. The benefit
grows with the share of human attackers ($+0.5\%$ against bot-dominated traffic, $+5.2\%$
against human-dominated traffic): adaptation pays only where attackers themselves adapt.

\vspace{0.5em}
\noindent\textbf{Keywords:} honeypot, SSH, deception, reinforcement learning, Q-learning,
Markov decision process, POMDP, stochastic approximation, event-driven architecture.
